\chapter{2012年安徽卷（文）}

\section{单选题}

\begin{problem}
复数 \(z\) 满足 \((z - \i)\i = 2 + \i\)，则 \(z =\)（\quad）

\choices
    {\(-1 - \i\)}
    {\(1 - \i\)}
    {\(-1 + 3\i\)}
    {\(1 - 2\i\)}
\end{problem}

\begin{answer}
B
\end{answer}

\begin{solution}
由 \((z-\i)\i=2+\i\)，两边同除以 \(\i\)，得 \(z-\i=\frac{2+\i}{\i}=1-2\i\)，所以 \(z=1-\i\). 故选 B.
\end{solution}

\begin{problem}
设集合 \(A = \{x \mid -3 \leqslant 2x - 1 \leqslant 3\}\)，集合 \(B\) 为函数 \(y = \lg (x - 1)\) 的定义域，则 \(A \cap B =\)（\quad）

\choices
    {\((1, 2)\)}
    {\([1, 2]\)}
    {\([1,2)\)}
    {\((1,2]\)}
\end{problem}

\begin{answer}
D
\end{answer}

\begin{solution}
由 \(-3\leqslant 2x-1\leqslant 3\)，得 \(-1\leqslant x\leqslant 2\)，所以 \(A=[-1,2]\). 函数 \(y=\lg(x-1)\) 要求 \(x-1>0\)，故 \(B=(1,+\infty)\). 因此 \(A\cap B=(1,2]\)，故选 D.
\end{solution}

\begin{problem}
\(\log_29\cdot \log_34 =\) （\quad）

\choices
    {\(\frac{1}{4}\)}
    {\(\frac{1}{2}\)}
    {2}
    {4}
\end{problem}

\begin{answer}
D
\end{answer}

\begin{solution}
利用换底公式，\(\log_2 9=2\log_2 3\)，\(\log_3 4=2\log_3 2\). 又 \(\log_2 3\cdot\log_3 2=1\)，所以 \(\log_2 9\cdot\log_3 4=4\)，故选 D.
\end{solution}

\begin{problem}
命题“存在实数 \(x\)，使 \(x > 1\)”的否定是（\quad）

\choices
    {对任意实数 \(x\)，都有 \(x > 1\)}
    {不存在实数 \(x\)，使 \(x \leqslant 1\)}
    {对任意实数 \(x\)，都有 \(x \leqslant 1\)}
    {存在实数 \(x\)，使 \(x \leqslant 1\)}
\end{problem}

\begin{answer}
C
\end{answer}

\begin{solution}
原命题是“存在 \(x\in\R\)，使 \(x>1\)”. 否定存在命题要改为对任意实数，并将不等式取否定，\(x>1\) 的否定为 \(x\leqslant 1\). 因此否定为“对任意实数 \(x\)，都有 \(x\leqslant 1\)”，故选 C.
\end{solution}

\begin{problem}
公比为 2 的等比数列 \(\{a_{n}\}\) 的各项都是正数，且 \(a_{3}a_{11} = 16\)，则 \(a_{5} =\)（\quad）

\choices
    {1}
    {2}
    {4}
    {8}
\end{problem}

\begin{answer}
A
\end{answer}

\begin{solution}
设等比数列公比为 \(q=2\). 由等比数列性质，\(a_3a_{11}=a_7^2=16\). 因各项为正数，故 \(a_7=4\). 又 \(a_7=a_5q^2=4a_5\)，所以 \(a_5=1\)，故选 A.
\end{solution}

\begin{problem}
如图所示，程序框图 （算法流程图） 的输出结果是（\quad）

\bitmapfigure[float=inline,max width=0.70\linewidth]{img/2012/anhui_liberal/q6_fig1.png}
\FloatBarrier

\choices
    {3}
    {4}
    {5}
    {8}
\end{problem}

\begin{answer}
B
\end{answer}

\begin{solution}
程序开始时 \(x=1\)，\(y=1\). 每次判断若 \(x\leqslant 4\)，就执行 \(x=2x\)，\(y=y+1\). 依次得到 \((x,y)=(2,2)\)、\((4,3)\)、\((8,4)\). 当 \(x=8>4\) 时停止循环，输出 \(y=4\)，故选 B.
\end{solution}

\begin{problem}
要得到函数 \( y = \cos (2x + 1) \) 的图象，只要将函数 \( y = \cos 2x \) 的图象（\quad）

\choices
    {向左平移 1 个单位}
    {向右平移 1 个单位}
    {向左平移 \(\frac{1}{2}\) 个单位}
    {向右平移 \(\frac{1}{2}\) 个单位}
\end{problem}

\begin{answer}
C
\end{answer}

\begin{solution}
令 \(g(x)=\cos 2x\)，则 \(g\left(x+\frac{1}{2}\right)=\cos(2x+1)\). 因此把 \(y=\cos 2x\) 的图象向左平移 \(\frac{1}{2}\) 个单位即可得到 \(y=\cos(2x+1)\) 的图象，故选 C.
\end{solution}

\begin{problem}
若 \(x\)，\(y\) 满足约束条件 \(\begin{cases}
x \geqslant 0,\\
x + 2y \geqslant 3,\\
2x + y \leqslant 3,
\end{cases}\) 则 \(z = x - y\) 的最小值是（\quad）

\choices
    {\(-3\)}
    {0}
    {\( \frac{3}{2} \)}
    {3}
\end{problem}

\begin{answer}
A
\end{answer}

\begin{solution}
可行域由 \(x\geqslant0\)、\(x+2y\geqslant3\)、\(2x+y\leqslant3\) 围成. 三条边界的相关交点为 \((0,\frac{3}{2})\)、\((0,3)\)、\((1,1)\)：其中后一个由
\(\begin{cases}
x+2y=3,\\
2x+y=3
\end{cases}\) 解得. 在这三个顶点处，\(z=x-y\) 的值分别为 \(-\frac{3}{2}\)、\(-3\)、\(0\). 线性目标函数在三角形可行域的顶点处取得最值，所以最小值为 \(-3\)，故选 A.
\end{solution}

\begin{problem}
若直线 \(x - y + 1 = 0\) 与圆 \((x - a)^2 + y^2 = 2\) 有公共点，则实数 \(a\) 的取值范围是（\quad）

\choices
    {\([-3, - 1]\)}
    {\([-1, 3]\)}
    {\([-3, 1]\)}
    {\((- \infty, -3] \cup [1, +\infty)\)}
\end{problem}

\begin{answer}
C
\end{answer}

\begin{solution}
圆心为 \((a,0)\)，半径为 \(\sqrt{2}\). 直线写成 \(x-y+1=0\)，圆心到直线的距离为 \(\frac{|a+1|}{\sqrt{2}}\). 直线与圆有公共点当且仅当该距离不大于半径，因此 \(\frac{|a+1|}{\sqrt{2}}\leqslant\sqrt{2}\)，即 \(|a+1|\leqslant2\). 解得 \(-3\leqslant a\leqslant1\)，故选 C.
\end{solution}

\begin{problem}
袋中共有 6 个除了颜色外完全相同的球，其中有 1 个红球，2 个白球和 3 个黑球. 从袋中任取两球，两球颜色为一白一黑的概率等于（\quad）

\choices
    {\( \frac{1}{5} \)}
    {\( \frac{2}{5} \)}
    {\( \frac{3}{5} \)}
    {\( \frac{4}{5} \)}
\end{problem}

\begin{answer}
B
\end{answer}

\begin{solution}
从 \(6\) 个球中任取两球的总数为 \(\mathrm C_6^2=15\). 取到一白一黑时，先从两个白球中取一个，再从三个黑球中取一个，共有 \(\mathrm C_2^1\mathrm C_3^1=6\) 种. 所求概率为 \(\frac{6}{15}=\frac{2}{5}\)，故选 B.
\end{solution}

\section{填空题}

\begin{problem}
设向量 \(\bs{a} = (1, 2m)\)，\(\bs{b} = (m + 1, 1)\)，\(\bs{c} = (2, m)\)，若 \((\bs{a} + \bs{c}) \perp \bs{b}\)，则 \(|\bs{a}| = \)\fillinblank{}.
\end{problem}

\begin{answer}
\(\sqrt{2}\)
\end{answer}

\begin{solution}
由垂直关系，\((\bs{a}+\bs{c})\cdot\bs{b}=0\). 因为 \(\bs{a}+\bs{c}=(3,3m)\)，\(\bs{b}=(m+1,1)\)，所以 \(3(m+1)+3m=0\)，解得 \(m=-\frac{1}{2}\). 此时 \(\bs{a}=(1,-1)\)，故 \(|\bs{a}|=\sqrt{1^2+(-1)^2}=\sqrt{2}\).
\end{solution}

\begin{problem}
某几何体的三视图如图所示，该几何体的体积是\fillinblank{}.

\bitmapfigure[float=inline,max width=0.58\linewidth]{img/2012/anhui_liberal/q12_fig1.png}
\FloatBarrier
\end{problem}

\begin{answer}
56
\end{answer}

\begin{solution}
由三视图可知，该几何体的底面是上、下底分别为 \(5\)、\(2\)，高为 \(4\) 的梯形，几何体的高为 \(4\). 因而底面积为 \(\frac{5+2}{2}\times4=14\)，体积为 \(14\times4=56\).
\end{solution}

\begin{problem}
若函数 \( f(x) = |2x + a| \) 的单调递增区间是 \([3, +\infty)\)，则 \( a = \)\fillinblank{}.
\end{problem}

\begin{answer}
\(-6\)
\end{answer}

\begin{solution}
函数 \(f(x)=|2x+a|\) 的折点满足 \(2x+a=0\)，即 \(x=-\frac{a}{2}\). 因为系数 \(2>0\)，函数在 \(\left[-\frac{a}{2},+\infty\right)\) 上单调递增. 由题意 \(-\frac{a}{2}=3\)，所以 \(a=-6\).
\end{solution}

\begin{problem}
过抛物线 \( y^{2}=4x \) 的焦点 \(F\) 的直线交该抛物线于 \(A\)，\(B\) 两点，若 \( |AF|=3 \)，则 \( |BF|= \)\fillinblank{}.
\end{problem}

\begin{answer}
\(\frac{3}{2}\)
\end{answer}

\begin{solution}
将抛物线 \(y^2=4x\) 参数化为 \(P(t)=(t^2,2t)\)，设 \(A=P(t_1)\)，\(B=P(t_2)\). 若 \(t_1+t_2=0\)，则直线 \(AB\) 为竖直直线 \(x=t_1^2\). 因为它过焦点 \(F=(1,0)\)，所以 \(t_1^2=1\)，从而 \(|AF|=t_1^2+1=2\)，与题设矛盾，故 \(t_1+t_2\ne0\). 于是过这两点的直线斜率为
\(\frac{2t_2-2t_1}{t_2^2-t_1^2}=\frac{2}{t_1+t_2}\)，所以直线方程为
\(y-2t_1=\frac{2}{t_1+t_2}(x-t_1^2)\)，整理得
\(2x-(t_1+t_2)y+2t_1t_2=0\). 抛物线焦点为 \(F=(1,0)\)，代入直线方程得 \(t_1t_2=-1\).

点 \(P(t)\) 到焦点的距离为 \(FP(t)=\sqrt{(t^2-1)^2+(2t)^2}=t^2+1\). 由 \(|AF|=3\) 得 \(t_1^2=2\)，从而 \(t_2^2=\frac{1}{2}\). 因此 \(|BF|=t_2^2+1=\frac{3}{2}\).
\end{solution}

\begin{problem}
若四面体 \(ABCD\) 的三组对棱分别相等，即 \(AB = CD\)，\(AC = BD\)，\(AD = BC\)，则\fillinblank{}.（写出所有正确结论编号）

\begin{circlelist}
\item 四面体 \(ABCD\) 每组对棱相互垂直；
\item 四面体 \(ABCD\) 每个面的面积相等；
\item 从四面体 \(ABCD\) 每个顶点出发的三条棱两两夹角之和大于 \(90^{\circ}\) 而小于 \(180^{\circ}\)；
\item 连接四面体 \(ABCD\) 每组对棱中点的线段相互垂直平分；
\item 从四面体 \(ABCD\) 每个顶点出发的三条棱的长可作为一个三角形的三边长.
\end{circlelist}
\end{problem}

\begin{answer}
\(\circled{2}\circled{4}\circled{5}\)
\end{answer}

\begin{solution}
设 \(\bs{u}=\overrightarrow{AB}\)，\(\bs{v}=\overrightarrow{AC}\)，\(\bs{w}=\overrightarrow{AD}\)，并记 \(p=|\bs{u}|\)、\(q=|\bs{v}|\)、\(r=|\bs{w}|\). 由三组对棱相等，有 \(|\bs{w}-\bs{v}|=p\)、\(|\bs{w}-\bs{u}|=q\)、\(|\bs{v}-\bs{u}|=r\). 展开平方后得到 \(\bs{u}\cdot\bs{v}=\frac{p^2+q^2-r^2}{2}\)、\(\bs{u}\cdot\bs{w}=\frac{p^2+r^2-q^2}{2}\)、\(\bs{v}\cdot\bs{w}=\frac{q^2+r^2-p^2}{2}\).

四面体的四个面的边长都分别为 \(p\)，\(q\)，\(r\)：例如面 \(ABC\) 的三边为 \(p\)，\(q\)，\(r\)，其余三个面也只是这三个长度的排列，所以四个面按边边边全等，面积相等，第 \(2\) 项结论正确.

设 \(AB\)、\(CD\) 的中点分别为 \(M_1\)、\(N_1\)，\(AC\)、\(BD\) 的中点分别为 \(M_2\)、\(N_2\)，\(AD\)、\(BC\) 的中点分别为 \(M_3\)、\(N_3\). 以 \(A\) 为原点，则三条中点连线的方向向量可分别取为
\(\bs{r}_1=\bs{v}+\bs{w}-\bs{u}\)、\(\bs{r}_2=\bs{u}+\bs{w}-\bs{v}\)、\(\bs{r}_3=\bs{u}+\bs{v}-\bs{w}\). 利用上面的数量积关系逐一计算：
\(\bs{r}_1\cdot\bs{r}_2=2\bs{u}\cdot\bs{v}+r^2-p^2-q^2=0\)，
\(\bs{r}_1\cdot\bs{r}_3=2\bs{u}\cdot\bs{w}+q^2-p^2-r^2=0\)，
\(\bs{r}_2\cdot\bs{r}_3=2\bs{v}\cdot\bs{w}+p^2-q^2-r^2=0\)，所以这三个向量两两垂直. 同时，\(M_1N_1\)、\(M_2N_2\)、\(M_3N_3\) 的中点都对应于向量 \(\frac{\bs{u}+\bs{v}+\bs{w}}{4}\)，所以三条中点连线相互垂直平分，第 \(4\) 项结论正确.

第 \(1\) 项结论不一定成立. 例如取 \(A=(1,2,1)\)、\(B=(1,-2,-1)\)、\(C=(-1,2,-1)\)、\(D=(-1,-2,1)\)，则
\(AB^2=CD^2=20\)、\(AC^2=BD^2=8\)、\(AD^2=BC^2=20\)，三组对棱分别相等；但
\(\overrightarrow{AB}\cdot\overrightarrow{CD}=(0,-4,-2)\cdot(0,-4,2)=12\ne0\)，所以对棱不一定垂直.

第 \(3\) 项结论不一定成立. 正四面体满足三组对棱分别相等，而每个顶点处的三个两两夹角均为 \(60^{\circ}\)，其和等于 \(180^{\circ}\)，不满足“小于 \(180^{\circ}\)”.

在每个顶点处，三条棱的长度都是 \(p\)，\(q\)，\(r\) 的排列. 由于它们是同一面的三边，例如面 \(ABC\) 给出 \(p<q+r\)、\(q<p+r\)、\(r<p+q\)，所以三条棱长可作为三角形的三边，第 \(5\) 项结论正确. 因此所有正确结论的编号为 \(\circled{2}\circled{4}\circled{5}\).
\end{solution}

\section{解答题}

\begin{problem}
设 \(\triangle ABC\) 的内角 \(A\)，\(B\)，\(C\) 所对的边为 \(a\)，\(b\)，\(c\)，且有 \(2\sin B\cos A = \sin A\cos C + \cos A\sin C\).

\begin{enumerate}
\item 求角 \(A\) 的大小；
\item 若 \(b = 2\)，\(c = 1\)，\(D\) 为 \(BC\) 的中点，求 \(AD\) 的长.
\end{enumerate}
\end{problem}

\begin{solution}
\begin{enumerate}
\item 由三角形内角和 \(A+B+C=\pi\)，得 \(A+C=\pi-B\). 因此题设等式右边为
\(\sin A\cos C+\cos A\sin C=\sin(A+C)=\sin(\pi-B)=\sin B\). 从而 \(2\sin B\cos A=\sin B\). 因为 \(0<B<\pi\)，所以 \(\sin B>0\)，可约去 \(\sin B\)，得 \(2\cos A=1\). 又 \(0<A<\pi\)，故 \(A=\frac{\pi}{3}\).
\item 由余弦定理，\(a^2=b^2+c^2-2bc\cos A=2^2+1^2-2\times2\times1\times\frac{1}{2}=3\). 因为 \(D\) 是 \(BC\) 的中点，利用中线长公式得 \(AD^2=\frac{2b^2+2c^2-a^2}{4}=\frac{2\times2^2+2\times1^2-3}{4}=\frac{7}{4}\). 所以 \(AD=\frac{\sqrt{7}}{2}\).
\end{enumerate}
\end{solution}

\begin{problem}
设 \(f(x) = ax + \frac{1}{ax} + b\) （\(a > 0\)） 定义在 \((0, +\infty)\) 上.

\begin{enumerate}
\item 求 \(f(x)\) 的最小值；
\item 若曲线 \(y = f(x)\) 在点 \((1, f(1))\) 处的切线方程为 \(y = \frac{3}{2}x\)，求 \(a\)，\(b\) 的值.
\end{enumerate}
\end{problem}

\begin{solution}
\begin{enumerate}
\item 因为 \(a>0\) 且 \(x>0\)，令 \(t=ax\)，则 \(t>0\)，并且 \(f(x)=t+\frac{1}{t}+b\). 由基本不等式，\(t+\frac{1}{t}\geqslant2\)，等号当且仅当 \(t=1\)，即 \(x=\frac{1}{a}\). 所以 \(f(x)\) 的最小值为 \(2+b\).
\item \(f(x)=ax+\frac{1}{ax}+b\) 的导数为 \(f'(x)=a-\frac{1}{ax^2}\). 切线斜率为 \(\frac{3}{2}\)，故在 \(x=1\) 处有 \(a-\frac{1}{a}=\frac{3}{2}\). 整理得 \(2a^2-3a-2=0\)，即 \((2a+1)(a-2)=0\). 因为 \(a>0\)，所以 \(a=2\).

切线方程为 \(y=\frac{3}{2}x\)，点 \((1,f(1))\) 在切线上，故 \(f(1)=\frac{3}{2}\). 代入 \(a=2\)，得 \(2+\frac{1}{2}+b=\frac{3}{2}\)，所以 \(b=-1\).
\end{enumerate}
\end{solution}

\begin{problem}
若某产品的直径长与标准值的差的绝对值不超过 \(1\mathrm{~mm}\) 时，则视为合格品，否则视为不合格品. 在近期一次产品抽样检查中，从某厂生产的此种产品中随机抽取 5000 件进行检测，结果发现有 50 件不合格品. 计算这 50 件不合格品的直径长与标准值的差（单位：mm），将所得数据分组，得到如下频率分布表：

\begin{center}
\begin{tabular}{c|c|c}
\hline
分组 & 频数 & 频率 \\
\hline
\([-3,-2)\) &\fillinblank{} & 0.10 \\
\([-2,-1)\) & 8 &\fillinblank{} \\
\((1,2]\) &\fillinblank{} & 0.50 \\
\((2,3]\) & 10 &\fillinblank{} \\
\((3,4]\) &\fillinblank{} &\fillinblank{} \\
\hline
合计 & 50 & 1.00 \\
\hline
\end{tabular}
\end{center}

\begin{enumerate}
\item 将上面表格中缺少的数据补充完整；
\item 估计该厂生产的这种产品中，不合格品的直径长与标准值的差落在区间 \((1,3]\) 内的概率；
\item 现对该厂这种产品的某个批次进行检查，结果发现有 20 件不合格品，据此估算这批产品中的合格品的件数.
\end{enumerate}
\end{problem}

\begin{solution}
\begin{enumerate}
\item 频率等于频数除以总频数，因此第一组频数为 \(50\times0.10=5\)，第二组频率为 \(\frac{8}{50}=0.16\)，第三组频数为 \(50\times0.50=25\)，第四组频率为 \(\frac{10}{50}=0.20\). 前四组频率之和为 \(0.10+0.16+0.50+0.20=0.96\)，所以第五组频率为 \(1-0.96=0.04\)，第五组频数为 \(50\times0.04=2\). 缺少的数据依次为 \(5\)，\(0.16\)，\(25\)，\(0.20\)，\(2\)，\(0.04\).
\item 差落在 \((1,3]\) 内的样本包括分组 \((1,2]\) 和 \((2,3]\)，共有 \(25+10=35\) 件. 因此估计所求概率为 \(\frac{35}{50}=0.70\).
\item 抽样中不合格品率为 \(\frac{50}{5000}=0.01\). 该批次有 \(20\) 件不合格品，估计该批产品总数为 \(\frac{20}{0.01}=2000\) 件，所以合格品约为 \(2000-20=1980\) 件.
\end{enumerate}
\end{solution}

\begin{problem}
如图，长方体 \(ABCD - A_{1}B_{1}C_{1}D_{1}\) 中，底面 \(A_{1}B_{1}C_{1}D_{1}\) 是正方形，\(O\) 是 \(BD\) 的中点，\(E\) 是棱 \(AA_{1}\) 上任意一点.

\begin{enumerate}
\item 证明： \(BD \perp EC_{1}\)；
\item 如果 \(AB = 2\)，\(AE = \sqrt{2}\)，\(OE \perp EC_1\)，求 \(AA_1\) 的长.
\end{enumerate}

\bitmapfigure[float=inline,max width=0.42\linewidth]{img/2012/anhui_liberal/q19_fig1.png}
\FloatBarrier
\end{problem}

\begin{solution}
\begin{enumerate}
\item 设正方形底面边长为 \(s\)，长方体高为 \(h\). 建立直角坐标系，取 \(A_1=(0,0,0)\)、\(B_1=(s,0,0)\)、\(D_1=(0,s,0)\)、\(C_1=(s,s,0)\)，并令 \(A=(0,0,h)\)、\(B=(s,0,h)\)、\(D=(0,s,h)\). 由于 \(O\) 是 \(BD\) 的中点，\(O=(\frac{s}{2},\frac{s}{2},h)\). 设 \(E=(0,0,t)\)，则 \(\overrightarrow{BD}=(-s,s,0)\)，\(\overrightarrow{EC_1}=(s,s,-t)\)，从而 \(\overrightarrow{BD}\cdot\overrightarrow{EC_1}=-s^2+s^2=0\). 所以 \(BD\perp EC_1\).
\item 由 \(AB=2\)，得 \(s=2\). 又 \(AE=\sqrt{2}\)，故 \(E=(0,0,h-\sqrt{2})\). 于是 \(\overrightarrow{OE}=(-1,-1,-\sqrt{2})\)，\(\overrightarrow{EC_1}=(2,2,-h+\sqrt{2})\). 条件 \(OE\perp EC_1\) 给出
\(\overrightarrow{OE}\cdot\overrightarrow{EC_1}=-2-2+\sqrt{2}(h-\sqrt{2})=0\).
因此 \(\sqrt{2}h-6=0\)，解得 \(AA_1=h=3\sqrt{2}\).
\end{enumerate}
\end{solution}

\begin{problem}
如图，\(F_{1}\)，\(F_{2}\) 分别是椭圆 \(C: \frac{x^{2}}{a^{2}} + \frac{y^{2}}{b^{2}} = 1 \) （\(a > b > 0\)）的左，右焦点，\(A\) 是椭圆 \(C\) 的顶点；\(B\) 是直线 \(AF_{2}\) 与椭圆 \(C\) 的另一个交点，\(\angle F_{1}AF_{2} = 60^{\circ}\).

\begin{enumerate}
\item 求椭圆 \(C\) 的离心率；
\item 已知 \(\triangle AF_{1}B\) 面积为 \(40\sqrt{3}\)，求 \(a\)，\(b\) 的值.
\end{enumerate}

\bitmapfigure[float=inline,max width=0.42\linewidth]{img/2012/anhui_liberal/q20_fig1.png}
\FloatBarrier
\end{problem}

\begin{solution}
\begin{enumerate}
\item 设椭圆焦距的一半为 \(c\)，则 \(F_1=(-c,0)\)、\(F_2=(c,0)\)，顶点 \(A=(0,b)\). 由 \(a^2=b^2+c^2\)，有 \(AF_1=AF_2=a\). 在等腰三角形 \(AF_1F_2\) 中，\(\angle F_1AF_2=60^{\circ}\)，由余弦定理得
\((2c)^2=a^2+a^2-2a^2\cos60^{\circ}=a^2\).
所以 \(c=\frac{a}{2}\)，椭圆离心率为 \(e=\frac{c}{a}=\frac{1}{2}\).
\item 直线 \(AF_2\) 的方程为 \(y=b-\frac{b}{c}x\). 将其代入椭圆方程，得
\(\frac{x^2}{a^2}+\left(1-\frac{x}{c}\right)^2=1\)，即
\(x\left[\left(\frac{1}{a^2}+\frac{1}{c^2}\right)x-\frac{2}{c}\right]=0\).
其中 \(x=0\) 对应交点 \(A\)，另一个交点 \(B\) 的横坐标为
\(x_B=\frac{2/c}{1/a^2+1/c^2}=\frac{2a^2c}{a^2+c^2}\). 代入 \(c=\frac{a}{2}\)，得 \(x_B=\frac{4a}{5}\)，进而 \(y_B=b\left(1-\frac{x_B}{c}\right)=-\frac{3b}{5}\). 因此 \(A=(0,b)\)、\(F_1=(-\frac{a}{2},0)\)、\(B=(\frac{4a}{5},-\frac{3b}{5})\).

三角形 \(AF_1B\) 的面积为
\(\frac{1}{2}\left|\det\begin{pmatrix}-\frac{a}{2}&-b\\ \frac{4a}{5}&-\frac{8b}{5}\end{pmatrix}\right|=\frac{4ab}{5}\).
由面积条件得 \(ab=50\sqrt{3}\). 又 \(b^2=a^2-c^2=a^2-\frac{a^2}{4}=\frac{3a^2}{4}\)，且 \(b>0\)，所以 \(b=\frac{\sqrt{3}}{2}a\). 于是 \(\frac{\sqrt{3}}{2}a^2=50\sqrt{3}\)，得 \(a=10\)，\(b=5\sqrt{3}\).
\end{enumerate}
\end{solution}

\begin{problem}
设 \(f(x) = \frac{x}{2} + \sin x\) 的所有正的极小值点从小到大排成的数列为 \(\{x_n\}\).

\begin{enumerate}
\item 求数列 \(\{x_n\}\) 的通项公式；
\item 设 \(\{x_n\}\) 的前 \(n\) 项和为 \(S_n\)，求 \(\sin S_n\).
\end{enumerate}
\end{problem}

\begin{solution}
\begin{enumerate}
\item \(f'(x)=\frac{1}{2}+\cos x\). 临界点满足 \(\cos x=-\frac{1}{2}\)，所以 \(x=\frac{2\pi}{3}+2k\pi\) 或 \(x=\frac{4\pi}{3}+2k\pi\)，其中 \(k\in\Z\). 又 \(f''(x)=-\sin x\)，在 \(x=\frac{2\pi}{3}+2k\pi\) 处 \(f''(x)=-\frac{\sqrt{3}}{2}<0\)，为极大值点；在 \(x=\frac{4\pi}{3}+2k\pi\) 处 \(f''(x)=\frac{\sqrt{3}}{2}>0\)，为极小值点. 正的极小值点从小到大为 \(\frac{4\pi}{3}\)，\(\frac{10\pi}{3}\)，\(\frac{16\pi}{3}\)，\(\ldots\)，故
\(x_n=\frac{4\pi}{3}+2(n-1)\pi=-\frac{2\pi}{3}+2n\pi\)，其中 \(n\in\mathbb N_{+}\).
\item \(S_n=\sum_{k=1}^n\left(-\frac{2\pi}{3}+2k\pi\right)=-\frac{2n\pi}{3}+\pi n(n+1)\). 因为 \(n(n+1)\) 为偶数，\(\pi n(n+1)\) 是 \(2\pi\) 的整数倍，所以 \(\sin S_n=\sin\left(-\frac{2n\pi}{3}\right)\). 进一步按 \(n\) 除以 \(3\) 的余数分类，\(\sin S_n=0\)（\(3\mid n\)），\(\sin S_n=-\frac{\sqrt{3}}{2}\)（\(n=3k+1\)），\(\sin S_n=\frac{\sqrt{3}}{2}\)（\(n=3k+2\)）.
\end{enumerate}
\end{solution}
